\section*{Capítulo \ref{ch:b2:huckel} --- Teoria de Hückel e sistemas conjugados}

\begin{solution}{exo:b2:huckel:1}
Níveis do alila ($n = 3$): $\alpha + 1.414\beta$, $\alpha$, $\alpha - 1.414\beta$. Cátion (2 elétrons):
$2\alpha + 2.828\beta$; radical (3): $3\alpha + 2.828\beta$; ânion (4): $4\alpha + 2.828\beta$. Os
elétrons adicionais vão para o nível não ligante e só mudam a parte em $\alpha$.
\end{solution}

\begin{solution}{exo:b2:huckel:2}
$2(\alpha + 1.618\beta) + 2(\alpha + 0.618\beta) = 4\alpha + 4.472\beta$.
\end{solution}

\begin{solution}{exo:b2:huckel:3}
Benzeno (6), ânion ciclopentadienila (6), tropílio (6): \omterm{def:b2:huckel:aromatic}{aromáticos}. Cátion
ciclopentadienila (4), ciclobutadieno (4), ciclo-octatetraeno plano (8):
\omterm{def:b2:huckel:aromatic}{antiaromáticos} se planos (o ciclo-octatetraeno real tem forma de banheira e
não é \omterm{def:b2:huckel:aromatic}{aromático}).
\end{solution}

\begin{solution}{exo:b2:huckel:4}
Um octógono no círculo: níveis $\alpha + 2\beta$, $\alpha + 1.414\beta$ (duas vezes), $\alpha$ (duas vezes),
$\alpha - 1.414\beta$ (duas vezes), $\alpha - 2\beta$. Oito elétrons: dois no mais baixo,
quatro no primeiro par, um em cada orbital do par não ligante.
\end{solution}

\begin{solution}{exo:b2:huckel:5}
$4.472\beta - 4\beta = 0.472\beta$, isto é, $0.472 \times 75 = \qty{35}{kJ/mol}$, contra cerca de
\qty{10}{kJ/mol} pelos dados de hidrogenação: o modelo de Hückel, com $\beta$ ajustado
ao benzeno, superestima a conjugação das cadeias abertas.
\end{solution}

\begin{solution}{exo:b2:huckel:6}
$p_{12} = p_{34} = 0.894$, $p_{23} = 0.447$: as ligações das extremidades C1--C2 e
C3--C4 são as mais curtas; a ligação central é mais longa, mas mais curta que
uma ligação simples.
\end{solution}

\begin{solution}{exo:b2:huckel:7}
Ânion, seis elétrons: $2 \times 2 + 4 \times 0.618$, $E_\pi = 6\alpha + 6.472\beta$, uma camada
fechada: \omterm{def:b2:huckel:aromatic}{aromático}. Cátion, quatro: $2 \times 2 + 2 \times 0.618$, $4\alpha + 5.236\beta$, com um
elétron em cada orbital do par degenerado: \omterm{def:b2:huckel:aromatic}{antiaromático}, muito instável.
\end{solution}

\begin{solution}{exo:b2:huckel:8}
Anel de sete membros: $\alpha + 2\beta$, $\alpha + 1.247\beta$ (duas vezes), depois níveis
antiligantes. Seis elétrons, como no cátion \ce{C7H7+}, preenchem exatamente os
níveis ligantes: um cátion \omterm{def:b2:huckel:aromatic}{aromático}. O composto se ioniza facilmente em
tropílio e brometo.
\end{solution}

\begin{solution}{exo:b2:huckel:9}
$1.802 + 1.247 + 0.445 - 0.445 - 1.247 - 1.802 = 0$: a soma das energias é $6\alpha$.
\end{solution}

\begin{solution}{exo:b2:huckel:10}
$n = 6$: $m_k = 2\cos(k\pi/7) = 1.802$, $1.247$, $0.445$, $-0.445$, $-1.247$, $-1.802$. Six
elétrons: $E_\pi = 6\alpha + 2(1.802 + 1.247 + 0.445)\beta = 6\alpha + 6.988\beta$;
deslocalização $0.988\beta$.
\end{solution}

\begin{solution}{exo:b2:huckel:11}
No quadrado, os dois elétrons de energia mais alta ficam um em cada orbital de
um par degenerado não ligante. Alongar duas ligações opostas levanta a
degenerescência: um orbital desce, o outro sobe, e os dois elétrons se
emparelham no mais baixo, o que abaixa a energia. A molécula se estabiliza
como um retângulo com duas ligações duplas localizadas; ela continua
\omterm{def:b2:huckel:aromatic}{antiaromática} e extremamente reativa.
\end{solution}

\begin{solution}{exo:b2:huckel:12}
Matriz (em unidades de $\beta$, com $\alpha$ como origem) $\begin{pmatrix}0 & 1\\ 1 & 1\end{pmatrix}$:
$m = 1.618$ e $-0.618$, $E = \alpha + 1.618\beta$ (ligante) e $\alpha - 0.618\beta$. Para o
\omterm{def:b2:diatomic-mos:bonding}{orbital ligante}, $c_X = 1.618c_C$; logo $c_C^2 = 0.276$, $c_X^2 = 0.724$; com dois
elétrons, $q_C = 0.55$, $q_X = 1.45$: a densidade $\pi$ é atraída para X.
\end{solution}

\begin{solution}{pb:b2:huckel:1}
\textbf{1.} $-123.1 - (-4.3) = \qty{-118.8}{kJ/mol}$.
\textbf{2.} $-123.1 - 104.6 = \qty{-227.7}{kJ/mol}$.
\textbf{3.} $-123.1 - 82.9 = \qty{-206.0}{kJ/mol}$.
\textbf{4.} $3 \times (-118.8) = \qty{-356.4}{kJ/mol}$.
\textbf{5.} $356.4 - 206.0 = \qty{150}{kJ/mol}$.
\textbf{6.} O dieno libera $2 \times 118.8 - 227.7 = \qty{10}{kJ/mol}$ a menos que duas
ligações duplas isoladas: a conjugação de uma cadeia aberta dá um pequeno
ganho; o anel fechado do benzeno, quinze vezes mais.
\textbf{7.} Uma matriz $6 \times 6$ com 1 entre cada átomo e seus dois vizinhos (1--2, 2--3,
\dots, 6--1) e 0 nos demais lugares.
\textbf{8.} $m = 2\cos(2\pi k/6)$: $2$, $1$, $1$, $-1$, $-1$, $-2$: $E = \alpha + 2\beta$, $\alpha + \beta$ (duas vezes),
$\alpha - \beta$ (duas vezes), $\alpha - 2\beta$.
\textbf{9.} Dois elétrons em $\alpha + 2\beta$, quatro no par $\alpha + \beta$.
\textbf{10.} $E_\pi = 6\alpha + 8\beta$.
\textbf{11.} $3(2\alpha + 2\beta) = 6\alpha + 6\beta$.
\textbf{12.} $2\beta$, isto é, $2|\beta|$ de estabilização.
\textbf{13.} $2 + 1 + 1 - 1 - 1 - 2 = 0$: os seis níveis somam $6\alpha$.
\textbf{14.} $2|\beta| = \qty{150}{kJ/mol}$: $|\beta| = \qty{75}{kJ/mol}$.
\textbf{15.} $0.472 \times 75 = \qty{35}{kJ/mol}$, contra \qty{10}{kJ/mol}: o modelo ajustado
ao benzeno superestima a cadeia aberta.
\textbf{16.} Os orbitais ocupados podem ser tomados como $\psi_1 = (1, 1, 1, 1, 1, 1)/\sqrt6$,
$\psi_2 = (2, 1, -1, -2, -1, 1)/\sqrt{12}$ e $\psi_3 = (0, 1, 1, 0, -1, -1)/2$. Para a ligação
1--2: $p_{12} = 2(\frac16 + \frac{2}{12} + 0) = \frac23$; por simetria, toda ligação tem $p = 2/3$.
\textbf{17.} Entre os das ligações das extremidades (0.894) e da ligação central
(0.447) do butadieno.
\textbf{18.} As seis ligações são equivalentes por simetria; seu comprimento,
\qty{139.7}{pm}, fica entre a ligação dupla do eteno (\qty{133.9}{pm}) e a
ligação simples do etano (\qty{153.6}{pm}).
\textbf{19.} $0.988|\beta| = \qty{74}{kJ/mol}$, metade da do benzeno: fechar o anel
dobra a \omterm{def:b2:huckel:delocalisation}{energia de deslocalização}.
\textbf{20.} $\alpha + 2\beta$, $\alpha + 1.414\beta$ (duas vezes), $\alpha$ (duas vezes), $\alpha - 1.414\beta$ (duas vezes),
$\alpha - 2\beta$.
\textbf{21.} Os dois últimos elétrons vão um para cada orbital do par não
ligante: uma camada aberta, \omterm{def:b2:huckel:aromatic}{antiaromática}.
\textbf{22.} Ela se dobra em forma de banheira: os orbitais p das ligações
duplas vizinhas deixam de ser paralelos, as ligações se alternam, e a molécula
se comporta como um polieno.
\textbf{23.} Quatro elétrons: o mesmo par preenchido pela metade, pior porque a
molécula não consegue evitar a planaridade; ela se distorce em um retângulo e
é extremamente reativa.
\textbf{24.} Um sistema conjugado monocíclico plano é \omterm{def:b2:huckel:aromatic}{aromático} com $4N + 2$
elétrons $\pi$ e \omterm{def:b2:huckel:aromatic}{antiaromático} com $4N$.
\textbf{25.} $\boldsymbol{|\beta| \approx \qty{75}{kJ/mol}}$.
\end{solution}
